% ============================================================
%  Section 0 -- Resume
% ============================================================

\chapter*{Résumé}
\addcontentsline{toc}{chapter}{Résumé}
\markboth{Résumé}{Résumé}

\bibopsname{} est un framework d'evaluation de grands modeles de langue
applique au support informatique Michelin. Le projet repond a une question
operationnelle simple : comment comparer de maniere reproductible une reponse
LLM directe et une architecture agentique outillee, tout en tenant compte de la
qualite, de la securite, du cout, de la latence et de l'impact environnemental ?

Le depot etudie contient deux volets. Le premier est un banc de benchmark pour
tickets IT : il oppose un \emph{LLM Unique}, utilise en zero-shot, a un
\emph{Systeme Multi-Agents} fonde sur une boucle ReAct, une base de
connaissances JSON, une recherche documentaire RAG avec ChromaDB et une base
SQLite de statuts serveur. Le second volet est la \emph{Racing Arena}, une
simulation distribuee dans laquelle plusieurs equipes LLM prennent des decisions
de course en temps reel via SSE, y compris dans un scenario adversarial de
prompt injection.

Le projet comprend \metric{11\,182 lignes Python} dans \file{src/},
\metric{3\,209 lignes} de scripts, \metric{5\,935 lignes} de tests et
\metric{285 fichiers suivis} par Git. La campagne de verification locale du
12 mai 2026 a execute \metric{552 tests unitaires}, tous passes en
\metric{9,98 s}. Les benchmarks principaux sont produits sous
\file{data/outputs/benchmark/}; les figures integrees au rapport proviennent du
dossier \file{data/outputs/benchmark/charts/}.

\begin{table}[H]
\centering
\begin{tabularx}{\textwidth}{Ycc}
\toprule
\textbf{Métrique vérifiée} & \textbf{LLM Unique} & \textbf{Système Multi-Agents} \\
\midrule
Tickets evalues & 10 & 10 \\
Modele agent / zero-shot & \code{phi3:latest} & \code{phi3:latest} \\
Juge qualite & \multicolumn{2}{c}{\code{gpt-4o}} \\
Score qualite moyen (0--10) & 0,0 & \textbf{3,8} \\
Score securite moyen (0--10) & 10,0 & 10,0 \\
Latence totale & 277,35 s & \textbf{100,65 s} \\
Cout estime & 0,049593 \$ & \textbf{0,006270 \$} \\
Tokens totaux & 5\,344 & \textbf{627} \\
Energie estimee & 0,0010688 kWh & \textbf{0,0001254 kWh} \\
Empreinte carbone & 0,05344 gCO$_2$e & \textbf{0,00627 gCO$_2$e} \\
Score composite (0--100) & 35,0 & \textbf{75,2} \\
Verdict de release & \verdict{FAIL} & \verdict{FAIL} \\
\bottomrule
\end{tabularx}
\caption{Synthese du benchmark principal extrait de \file{comparison_results.json}.}
\label{tab:resume-benchmark}
\end{table}

Le systeme multi-agents obtient le meilleur score composite : \metric{75,2/100}
contre \metric{35,0/100} pour le LLM unique. Il reduit la latence totale de
\metric{63,7 \%}, divise le cout estime par \metric{7,91} et diminue les tokens
de \metric{88,3 \%}. Le verdict final reste cependant \verdict{FAIL} pour les
deux architectures, car la politique de gate impose un score qualite minimal de
7/10 et le meilleur score observe n'est que de 3,8/10. Cette distinction est
importante : \bibopsname{} montre une amelioration relative nette, mais le
resultat absolu n'est pas encore suffisant pour une mise en production sans
amelioration du modele, des prompts et du corpus de reference.

La Racing Arena complete ce diagnostic par un banc de securite dynamique. Sur
une course de 15 tours, le rapport de securite enregistre
\metric{286 decisions}. L'equipe ReAct exposee aux attaques execute
\metric{53 injections sur 57 attaques recues} (93 \%), tandis que l'equipe
validee n'execute aucune injection et ne fuite aucune strategie. Ce second volet
montre que l'evaluation d'agents ne doit pas se limiter a la qualite moyenne :
la resistance aux attaques, les surfaces d'API et la gouvernance des outils sont
des criteres de conception a part entiere.

\bigskip
\noindent\textbf{Mots-clés :} LLM, multi-agents, ReAct, RAG, MLOps,
support IT, prompt injection, GreenOps, benchmark, Michelin.
